\documentclass[tikz,border=6pt]{standalone}
% 공통 프리앰블 — PM Day1 교재 그림 (TikZ standalone)
\usepackage{fontspec}
\setmainfont{Apple SD Gothic Neo}
\setsansfont{Apple SD Gothic Neo}
\setmonofont{Menlo}
\usepackage{amssymb}
\usepackage{tikz}
\usetikzlibrary{arrows.meta,positioning,calc,shapes.geometric,fit,backgrounds,matrix,decorations.pathreplacing}
\definecolor{navy}{HTML}{1F3A5F}
\definecolor{teal}{HTML}{2A7F62}
\definecolor{rust}{HTML}{C8553D}
\definecolor{sand}{HTML}{E8E1D5}
\definecolor{ink}{HTML}{222222}
\definecolor{grey}{HTML}{7A7A7A}
\definecolor{lightgrey}{HTML}{EFEFEF}
\definecolor{navylight}{HTML}{DCE6F2}
\definecolor{teallight}{HTML}{DDF0E8}
\definecolor{rustlight}{HTML}{F6E0DA}
\tikzset{
  every node/.style={font=\small, text=ink},
  box/.style={draw=navy, thick, rounded corners=2pt, align=center, inner sep=5pt, fill=white},
  boxfill/.style={box, fill=navylight},
  boxteal/.style={draw=teal, thick, rounded corners=2pt, align=center, inner sep=5pt, fill=teallight},
  boxrust/.style={draw=rust, thick, rounded corners=2pt, align=center, inner sep=5pt, fill=rustlight},
  boxgrey/.style={draw=grey, thick, rounded corners=2pt, align=center, inner sep=5pt, fill=lightgrey},
  arr/.style={-{Stealth[length=2.5mm]}, thick, navy},
  arrteal/.style={-{Stealth[length=2.5mm]}, thick, teal},
  arrrust/.style={-{Stealth[length=2.5mm]}, thick, rust},
  arrgrey/.style={-{Stealth[length=2.5mm]}, thick, grey},
  lbl/.style={font=\footnotesize, text=grey, align=center},
  src/.style={font=\scriptsize, text=grey, align=left},
  title/.style={font=\normalsize\bfseries, text=navy, align=center},
}

\begin{document}
\begin{tikzpicture}[node distance=6mm]
\node[title] (t) at (0,0) {두 개의 별개 관찰 — iron law와 곱셈 구조};
% 관찰 1
\node[boxrust, text width=64mm, below=6mm of t, xshift=-40mm] (law) {\textbf{관찰 1 — iron law of megaprojects}\\[2pt] 대형 프로젝트 열 중 아홉은 원가를 초과한다\\{\footnotesize "Over budget, over time, over and over again"}};
% 관찰 2
\node[title, font=\small\bfseries, below=6mm of t, xshift=40mm, text=navy] (t2) {관찰 2 — 세 차원의 곱셈 구조};
\node[boxfill, text width=17mm, below=4mm of t2, xshift=-24mm] (c) {원가\\{\footnotesize on budget}};
\node[boxfill, text width=17mm, right=4mm of c] (s) {일정\\{\footnotesize on time}};
\node[boxfill, text width=17mm, right=4mm of s] (b) {편익\\{\footnotesize on benefits}};
\node[font=\large, text=navy] at ($(c.east)!0.5!(s.west)$) {$\times$};
\node[font=\large, text=navy] at ($(s.east)!0.5!(b.west)$) {$\times$};
\node[lbl, below=2mm of c] {1/10}; \node[lbl, below=2mm of s] {1/10}; \node[lbl, below=2mm of b] {1/10};
\node[box, text width=58mm, below=10mm of s, fill=sand, draw=grey] (r1) {각 1/10이면 셋 동시 만족 $\approx$ \textbf{1/1,000}\\관대하게 각 2/10이어도 $\approx$ \textbf{8/1,000}};
% 하단 함의
\node[boxteal, text width=150mm, below=8mm of r1.south, anchor=north, xshift=-40mm] (imp) {\textbf{함의} — 조직은 대개 한 차원(대개 일정)만 관리하면서 프로젝트를 성공이라 부른다.\\온담 2019년 스마트오더: 34억을 예산 안에서 썼으나 8개월 만에 중단 → 원가 차원 성공, 편익 차원 실패.};
\draw[arrgrey] (r1.south) -- ($(r1.south)+(0,-8mm)$);
\node[src, below=3mm of imp.south west, anchor=north west] {출처: Flyvbjerg (2014), What You Should Know About Megaprojects and Why, Project Management Journal 45(2)를 바탕으로 재구성.};
\end{tikzpicture}
\end{document}
